%% Thanks to Bennet Goeckner for giving me his TeX template, which this is based on.
\documentclass[12pt]{article}

\usepackage{amsmath}
\usepackage{amsthm}
\usepackage{amsfonts}
\usepackage{amssymb}
\usepackage{enumerate}
\usepackage{graphicx}
\usepackage{bbm}
\usepackage{mdframed}
\usepackage{multicol}
\usepackage{verbatim}
\usepackage{tikz}
\usepackage{eurosym}
\usepackage{comment}
\usepackage{url}
\usepackage[margin = .8in]{geometry}
\usepackage[spanish, es-tabla]{babel}
\geometry{letterpaper}
\linespread{1.2}

\newcommand{\RR}{\mathbb{R}}
\newcommand{\NN}{\mathbb{N}}
\newcommand{\ZZ}{\mathbb{Z}}
\newcommand{\QQ}{\mathbb{Q}}

\newcommand\set[1]{\left\lbrace #1 \right\rbrace}
\newcommand\abs[1]{\left| #1 \right|}
\newcommand\parens[1]{\left( #1 \right)}
\newcommand\brac[1]{\left[ #1 \right]}
\newcommand\sol[1]{\begin{mdframed}
\emph{Solución.} #1
\end{mdframed}}
\newcommand\solproof[1]{\begin{mdframed}
\begin{proof} #1
\end{proof}
\end{mdframed}}

\begin{document}
\noindent Javier González Otero \hfill 01/06/2026

\begin{center}
    \vspace{1em}
    {\large \textbf{Examen Teórico}} \\
    {Fundamentos de Machine Learning}
    \vspace{1em}
\end{center}

\begin{enumerate}

\item \textbf{Interpretabilidad frente a capacidad predictiva.}

Un gestor de cartera afirma que no utilizará ningún modelo que no pueda explicar a sus clientes
institucionales. Sin embargo, el equipo de ciencia de datos sostiene que los modelos más complejos
ofrecen un rendimiento predictivo superior. Discute el compromiso existente entre interpretabilidad
y capacidad predictiva. Compara, desde esta perspectiva, modelos como Random Forest, GBM y
XGBoost. Explica, además, cómo las herramientas de interpretabilidad y las técnicas de selección de
variables pueden ayudar a comprender las decisiones de modelos complejos y facilitar su utilización
en entornos financieros.

\sol{

El compromiso entre interpretabilidad y capacidad predictiva surge de que los modelos con mayor
flexibilidad para capturar relaciones no lineales y de alta dimensión tienden a ser más difíciles de
explicar. Conviene matizar, no obstante, que este compromiso no es estrictamente binario: existen
técnicas que permiten acercarse a ambos objetivos de forma simultánea, de modo que la disyuntiva
entre la postura del gestor y la del equipo de ciencia de datos es menos tajante de lo que parece.

En cuanto a los modelos citados, el Random Forest es un ensamblaje de árboles entrenados mediante
bagging y selección aleatoria de variables; ofrece buena capacidad predictiva y cierta resistencia
al sobreajuste, es robusto frente a la elección de hiperparámetros y resulta fácil de paralelizar.
Su interpretabilidad nativa es limitada, pues no existe un único árbol que pueda examinarse, aunque
proporciona importancias de variables por reducción de impureza. El GBM (Gradient Boosting Machine),
por su parte, construye los árboles de forma secuencial, corrigiendo cada uno el error residual del
anterior; suele superar al Random Forest en datos tabulares estructurados, pero exige una afinación
más cuidadosa de los hiperparámetros (tasa de aprendizaje, profundidad, submuestreo) y es más
susceptible al sobreajuste si no se regula adecuadamente. XGBoost es una implementación optimizada
de gradient boosting que incorpora regularización $L_1$ y $L_2$, manejo nativo de valores faltantes
y una notable eficiencia computacional; se ha convertido en una de las opciones de referencia para
datos tabulares, y aunque su interpretabilidad nativa es comparable a la de GBM, se beneficia
especialmente de las herramientas externas de explicabilidad.

Son precisamente esas herramientas de interpretabilidad las que permiten reconciliar ambos extremos.
La importancia de variables por reducción de impureza, nativa de los árboles, da una primera lectura
global de qué variables pesan más, aunque puede sesgarse hacia las de alta cardinalidad. SHAP
(SHapley Additive exPlanations) afina esa lectura: descompone la predicción de cada observación en
contribuciones aditivas de las variables, y proporciona tanto explicaciones globales (importancia
media) como locales (un caso concreto), lo que resulta valioso para justificar decisiones
individuales ante clientes institucionales. LIME aproxima el comportamiento del modelo en el entorno
de una observación mediante un modelo lineal sencillo; es útil para comunicar decisiones puntuales,
si bien sus explicaciones pueden ser inestables. A todo ello se suma la selección de variables, que
al descartar las irrelevantes reduce la superficie que el analista debe explicar y hace más
transparente el modelo resultante.

En un entorno financiero, la combinación de XGBoost o GBM con SHAP permite atender la exigencia de
explicabilidad sin renunciar al rendimiento: el gestor puede mostrar, para cada señal de inversión,
qué variables la impulsan y en qué magnitud, y trasladar esa explicación al lenguaje del negocio (por
ejemplo, que el modelo sobrepondera una acción por la mejora de márgenes y la baja deuda neta).
}

%% ─────────────────────────────────────────────
\item \textbf{Diagnóstico y tratamiento del sobreajuste.}

Describe las señales prácticas que indican sobreajuste durante el desarrollo de un modelo y las
acciones correctivas que aplicaría, en orden de prioridad.

\sol{

Existen varias señales prácticas que delatan el sobreajuste durante el desarrollo. La más
característica es una brecha amplia entre la métrica de entrenamiento y la de validación o test
(por ejemplo, un AUC de 0{,}97 en entrenamiento frente a 0{,}72 en test); esa brecha se hace más
evidente cuando el error de entrenamiento sigue disminuyendo mientras el de validación se estanca o
empeora. También resulta sospechoso que el modelo se comporte de forma errática ante pequeñas
perturbaciones de los datos de entrada, que las importancias de variables revelen un peso exagerado
de variables ruidosas o de alta cardinalidad, o que la calibración sea deficiente, es decir, que las
probabilidades predichas no se correspondan con las frecuencias observadas en datos nuevos.

\medskip
En cuanto a las acciones correctivas, conviene aplicarlas en un orden que priorice las más robustas:

\begin{enumerate}[1.]
  \item Obtener más datos es una de las soluciones más robustas y preferibles cuando es viable, pues el
  sobreajuste es esencialmente un problema de muestra insuficiente en relación con la complejidad del
  modelo.
  \item Reducir la complejidad del modelo (menor profundidad de los árboles, menos estimadores o
  menos capas y neuronas) es la medida más directa cuando no se dispone de más datos.
  \item Aplicar regularización: penalizaciones $L_1$ o $L_2$ en modelos lineales y en XGBoost,
  dropout en redes neuronales, o restringir el tamaño de las hojas y las variables candidatas en los
  árboles.
  \item Recurrir a early stopping, monitorizando la métrica de validación durante el entrenamiento y
  deteniéndolo cuando deja de mejorar, algo especialmente útil en GBM y en redes neuronales.
  \item Seleccionar variables, ya que eliminar las irrelevantes o ruidosas reduce la superficie de
  memorización del modelo.
  \item Emplear validación cruzada estratificada, que proporciona una estimación más fiable del
  rendimiento real y permite distinguir si el sobreajuste es consistente o puntual.
  \item Como último recurso, generar datos adicionales o inyectar ruido controlado, que puede actuar
  como regularizador implícito.
\end{enumerate}
}

%% ─────────────────────────────────────────────
\item \textbf{Factores no técnicos que determinan el impacto de un modelo.}

Un modelo con métricas superiores a todo lo existente en una organización lleva seis meses sin ser
adoptado en la toma de decisiones. Analice las razones por las que esto ocurre con frecuencia y qué
factores, más allá de la calidad técnica, determinan que un modelo tenga un impacto real.

\sol{

Un modelo técnicamente excelente puede resultar irrelevante si no se integra en el flujo operativo
de quienes deben utilizarlo, y esto ocurre por varias razones.

La primera es la falta de confianza: los usuarios finales no comprenden cómo funciona el modelo y se
muestran reacios a decidir sobre la base de un sistema opaco. Generar esa confianza exige
explicabilidad, un historial observado de aciertos y, en muchas organizaciones, el respaldo de una
figura de autoridad interna. A ello se une la desalineación con el proceso de negocio: el modelo
puede producir predicciones en un formato, un momento o una granularidad que no encajan con la
decisión real, de modo que si el gestor decide cada semana y el modelo emite señales diarias sin
integrarse en su herramienta habitual, el coste de adopción supera al beneficio percibido.

Influye también la ausencia de un responsable interno con poder de decisión que respalde activamente
el modelo, sin el cual los incentivos para cambiar el statu quo resultan insuficientes. A esto se
añade la inercia organizativa y política, pues los modelos amenazan procesos establecidos o el
estatus de equipos que antes decidían de forma manual, de manera que la resistencia no es técnica
sino de poder. Otro factor recurrente es la falta de validación operativa: cuando las métricas de
desarrollo (AUC, F1) no se traducen a unidades de negocio comprensibles (rentabilidad incremental,
ahorro de costes, reducción de errores), la dirección carece de argumentos para justificar el
cambio. Finalmente, pesan la fiabilidad y la infraestructura, ya que un modelo que falla a menudo,
produce predicciones tardías o requiere intervención manual se abandona con rapidez, así como los
factores regulatorios y de gobernanza, que en entornos financieros pueden bloquear la adopción por
falta de documentación de auditoría, de validación por el equipo de riesgos o de aprobación del
comité de modelos.
}

%% ─────────────────────────────────────────────
\item \textbf{Construcción de un pipeline cuantitativo completo.}

Imagina que debes desarrollar un sistema para seleccionar acciones que formará parte del proceso de
inversión de un Hedge Fund. Dispones de cientos de variables financieras y de mercado, pero el tiempo
de cálculo y la interpretabilidad son factores importantes. Describe un pipeline completo que combine
reducción dimensional, selección de variables y modelos de ensamblaje. Justifica el orden de las
etapas, las técnicas que utilizarías en cada una de ellas y cómo controlarías el riesgo de
sobreajuste. Finalmente, explica qué métricas utilizarías para evaluar tanto la calidad estadística
del modelo como su utilidad económica en una estrategia de inversión real.

\sol{

Antes de cualquier modelado conviene fijar el diseño del problema y prevenir el data leakage: definir
el objetivo (por ejemplo, el retorno ajustado al riesgo a un horizonte de 21 días), establecer las
ventanas de entrenamiento y de validación walk-forward, y asegurarse de que ninguna variable
incorpora información posterior a la fecha de predicción. Esta etapa precede a todas las demás porque
un error en la definición del objetivo o una fuga de información invalidan el resto del pipeline por
bueno que sea.

A continuación se realiza el preprocesamiento y la limpieza: tratamiento de valores faltantes
mediante imputación por la mediana, robusta frente a outliers, o por media móvil en series
temporales; detección y tratamiento de valores atípicos; y estandarización cross-sectional (por
fecha) para neutralizar los efectos de escala y de mercado.

La reducción dimensional va antes que la selección de variables porque, con cientos de
variables, muchas estarán correlacionadas y conviene depurar primero esa redundancia.
Se eliminan las variables de baja varianza o muy correlacionadas ($|\rho| > 0{,}95$) y, a
continuación, se aplica PCA para resumir la redundancia restante, reteniendo las componentes que
expliquen entre el 80 y el 90\,\% de la varianza. Como filtro previo, conviene además conservar las
variables originales que tienen un peso apreciable en esas componentes y descartar las que no
contribuyen a ninguna, de modo que la reducción no solo comprima la información sino que también
ayude a depurar el conjunto de variables.

La selección de variables se aplica luego sobre las variables originales que han superado ese filtro,
en dos pasos. Primero un filtrado univariante mediante la correlación de Spearman con el objetivo o
la información mutua, que descarta las variables sin señal. Después una selección por importancia de
modelo, empleando un XGBoost o Random Forest ligero y regularizado del que se lee la importancia de
cada variable (por reducción de impureza o con SHAP), reteniendo las variables con mayor
contribución estable a lo largo del tiempo.

Sobre ese conjunto se entrenan los modelos base del ensamblaje, buscando que sean complementarios:
modelos lineales Ridge o Lasso como referencia de baja varianza, XGBoost para capturar no
linealidades e interacciones, LightGBM por su eficiencia de cálculo cuando hay que reentrenar con
frecuencia, y un modelo de factores lineales interpretable y alineado con la lógica financiera. Sus
predicciones se combinan mediante un promedio ponderado, asignando más peso a los modelos con mejor
rendimiento reciente para adaptarse a los cambios de régimen.

El control del sobreajuste descansa en tres pilares: una validación walk-forward, dejando una
separación temporal entre el entrenamiento y la validación para evitar que se solapen observaciones;
la regularización en todos los modelos base; y la monitorización fuera de muestra de la capacidad
predictiva a lo largo del tiempo, reiniciando el pipeline si decae.

Para la evaluación conviene combinar métricas estadísticas y económicas. Entre las primeras, la
correlación de Spearman entre la señal y el retorno realizado y la calidad de la calibración de las
probabilidades. Entre las económicas, el Sharpe ratio de la estrategia long/short, el máximo
drawdown y el turnover con su coste de transacción estimado, que indican si el acierto estadístico
llega a traducirse en valor para la estrategia.
}

%% ─────────────────────────────────────────────
\item \textbf{Generalización frente a memorización: una explicación rigurosa y accesible.}

Un interlocutor no especialista le pregunta: ``¿El modelo realmente aprende o simplemente
memoriza?'' Construya una respuesta que sea simultáneamente precisa y accesible para un profesional
inteligente sin formación técnica en ML.

\sol{

Ambas cosas pueden ocurrir, y saber diferenciarlas es lo más importante al crear estos sistemas.

\medskip
Un sistema memoriza cuando da respuestas perfectas para la información exacta que se usó para crearlo,
pero se equivoca en cuanto se le presenta información nueva. Esto ocurre porque el sistema se ha fijado
en detalles irrelevantes o casuales de esos datos iniciales, es decir, no es capaz de abstraer sus respuestas a un modelo general que cubra todos los casos posibles.

\medskip
En contraparte, podemos considerar que el sistema ''aprende'' cuando logra extraer e identificar las reglas o
características comunes que siempre se cumplen, es decir, logra abstraer.

\medskip
La forma de comprobar qué está haciendo el sistema es muy directa. Primero, el sistema se prepara con cierta información relevante al problema. Posteriormente, se le dan datos completamente nuevos que aporten información adicional, eventos no directamente imputables desde los datos iniciales... .
Si funciona bien tanto con los datos iniciales como con los nuevos, está abstrayendo y por ende podemos considerar que está aprendiendo; si solo
funciona bien con los datos iniciales y falla en los nuevos, simplemente ha memorizado.

\medskip
El objetivo de la persona que diseña el sistema es ajustarlo para encontrar un punto medio estricto: el
programa no debe ser tan básico que pase por alto las reglas importantes, ni tan rebuscado que acabe memorizando
detalles que no importan.

\medskip
Para asegurar que estamos ante un aprendizaje real y no ante pura memorización, el método principal consiste en
ocultar una parte de la información disponible desde el principio. Una parte de los datos se usa para construir
el sistema, y la parte que se mantuvo oculta se guarda exclusivamente para examinarlo al final y verificar que sus
respuestas siguen siendo correctas en un escenario nuevo.
}

%% ─────────────────────────────────────────────
\item \textbf{La definición del target como decisión de diseño.}

Dos equipos modelan el mismo fenómeno de abandono de clientes con definiciones diferentes: el equipo~A
define el abandono como ``ausencia de actividad en 90 días'', mientras que el equipo~B lo define como
``cancelación explícita de la suscripción''. Ambos obtienen métricas excelentes. Explique por qué uno
de estos modelos podría ser inútil, a pesar de su buen rendimiento técnico.

\sol{

La definición del objetivo no es un detalle técnico, sino la hipótesis de negocio que el modelo
valida. Si esa hipótesis no se alinea con la acción que la organización puede ejecutar, un modelo
perfectamente calibrado resulta inútil.

\medskip
El modelo del equipo~A, basado en la inactividad durante 90 días, probablemente predice con
precisión a clientes que ya están inactivos o en proceso de abandono silencioso. El problema es doble.
Por un lado, la inactividad puede ser temporal, estacional o irrelevante para el negocio, como en un
cliente que no actúa pero sigue pagando. Por otro, y más importante, si el modelo etiqueta a alguien
como abandono solo cuando ya lleva semanas sin actividad, la intervención comercial llega demasiado
tarde: el modelo describe el pasado en lugar de anticipar un futuro sobre el que aún se pueda actuar.

\medskip
El modelo del equipo~B, basado en la cancelación explícita, se ajusta en cambio al evento de negocio
que se quiere evitar. Es un evento claro, verificable y accionable: cuando el modelo predice que un
cliente va a cancelar, existe una ventana temporal para intervenir con una oferta de retención.

\medskip
El riesgo específico del equipo~A es que, aun con métricas excelentes, podría estar aprendiendo a
identificar patrones posteriores al hecho (clientes que ya abandonaron) en lugar de señales
predictivas tempranas; la variable ``días sin actividad'' puede actuar como un sustituto del propio
objetivo, lo que inflaría artificialmente el rendimiento sin aportar valor predictivo real.
}

%% ─────────────────────────────────────────────
\item \textbf{Segmentación de activos financieros y construcción de factores.}

Dispones de una base de datos con 2.000 acciones globales y más de 100 variables financieras,
fundamentales y de mercado para cada una de ellas. El objetivo es identificar grupos homogéneos de
compañías para construir estrategias de inversión diferenciadas. Explica cómo abordarías este
problema utilizando técnicas de reducción dimensional y clustering. Compara las ventajas e
inconvenientes de utilizar K-Means, DBSCAN y Clustering Jerárquico en este contexto. Discute,
además, qué métricas emplearías para evaluar la calidad de los clústeres obtenidos y cómo
interpretarías financieramente los resultados.

\sol{

Con más de cien variables, tenemos un espacio de alta dimensión, por lo que podemos tener curse of dimensionality:
las distancias euclídeas pierden significado y la mayoría de los algoritmos de clustering degradan su
rendimiento. Por ello es obligatorio normalizar las variables (estandarización cross-sectional) y
reducir la dimensionalidad antes de agrupar.

\medskip
Para la reducción dimensional, el PCA es rápido y reproducible y permite retener las componentes que
expliquen entre el 80 y el 85\,\% de la varianza, con el inconveniente de que esas componentes no son
directamente interpretables. Como alternativas no lineales, t-SNE y UMAP preservan mejor la
estructura local y resultan útiles para visualizar en dos dimensiones y detectar agrupaciones de
forma irregular, aunque son menos reproducibles y más difíciles de interpretar.

\medskip
En cuanto a los algoritmos de clustering:

\begin{itemize}
  \item K-Means es rápido, escala sin problema a 2.000 acciones y resulta fácil de interpretar, pues
  los centroides representan el perfil medio de cada grupo y el número de clústeres $k$ es un parámetro
  de diseño alineado con la granularidad deseada de la estrategia. En contra, asume clústeres
  esféricos y de tamaño similar, es sensible a los outliers y obliga a fijar $k$ de antemano, de modo
  que con distribuciones financieras asimétricas puede producir grupos heterogéneos.

  \item DBSCAN no exige fijar el número de clústeres, detecta formas arbitrarias e identifica
  automáticamente outliers, lo que puede ser
  valioso para la gestión de riesgos. Sus inconvenientes son la sensibilidad a los parámetros
  $\varepsilon$ y \textit{minPts}, su mal comportamiento cuando las densidades de los grupos difieren
  mucho y la degradación de la noción de densidad local en alta dimensión, incluso tras aplicar PCA.

  \item El clustering jerárquico tampoco requiere fijar $k$ y, mediante el dendrograma, permite
  explorar la jerarquía de fusiones, lo que es especialmente útil para construir estrategias a
  distintos niveles de granularidad (sectores frente a subsectores); el método de Ward es robusto y
  minimiza la varianza intragrupo. Su principal limitación es la complejidad computacional, de orden
  $O(n^2)$ a $O(n^3)$, que puede resultar prohibitiva para $n > 2.000$ salvo que se use una versión
  aproximada, y el hecho de que, una vez unidas dos observaciones, la decisión es irreversible.
\end{itemize}

\medskip
La calidad de los clústeres se evalúa con varias métricas complementarias: el Silhouette Score, que
mide la cohesión interna y la separación entre grupos (valores próximos a 1 indican clústeres bien
definidos); el índice de Davies-Bouldin, razón entre la dispersión intragrupo y la distancia entre
centroides, preferible cuanto menor sea; y, sobre todo, la estabilidad, que se comprueba repitiendo
el clustering sobre submuestras o con ruido y midiendo la consistencia de las asignaciones, pues un
agrupamiento útil debe ser estable.

\medskip
Finalmente, la interpretación financiera exige validar los clústeres con lógica de dominio:
conviene comprobar si los grupos se ordenan de forma natural por sector, capitalización, región o
perfil de crecimiento frente a valor. Si un clúster reúne tecnológicas de alto crecimiento y baja
rentabilidad y otro agrupa utilities de alta rentabilidad por dividendo, el agrupamiento tiene
sentido económico accionable. La validación más rigurosa consiste en comprobar si la rentabilidad
futura promedio difiere de forma apreciable entre clústeres, o si las estrategias long/short basadas
en ellos obtienen una rentabilidad superior ajustada al riesgo.
}

%% ─────────────────────────────────────────────
\item \textbf{Estrategia de modelado bajo restricción temporal.}

Dispone de un dataset tabular de tamaño moderado (50.000 observaciones, 30 variables), un problema
de clasificación binaria y un plazo de 24 horas. Describe su estrategia completa, justificando cada
decisión.

\sol{

Con un plazo de 24 horas y un dataset de tamaño moderado, la prioridad es maximizar el valor esperado
del resultado gestionando de forma realista el esfuerzo humano y computacional. La estrategia se organiza
en bloques lógicos, alternando trabajo activo y tiempo de cómputo desatendido.

\medskip
Las dos primeras horas se dedican al análisis exploratorio y a una limpieza rápida: revisar
distribuciones, valores faltantes, cardinalidad de las variables categóricas y balance de clases;
identificar posibles fugas de información (variables con una correlación sospechosamente alta con el
objetivo); y decidir si el desbalance requiere tratamiento, planificándolo si la clase positiva
representa menos del 5\,\%.

\medskip
Entre las horas dos y cuatro conviene construir un modelo de referencia sencillo (una regresión
logística o un LightGBM con parámetros por defecto y validación estratificada en 5 particiones).
Su rendimiento fija el umbral mínimo y permite extraer las primeras métricas de importancia para
guiar los siguientes pasos.

\medskip
Las horas cuatro a ocho se destinan a un \textit{feature engineering} selectivo basado en las
importancias (SHAP) del modelo de referencia: crear interacciones entre las tres a cinco variables
más relevantes con lógica de dominio, transformar variables asimétricas y eliminar las de importancia
nula. Realizar esto antes de la optimización garantiza que los hiperparámetros finales se ajusten
al conjunto de datos definitivo.

\medskip
Entre las horas ocho y dieciséis se configura y ejecuta el modelado principal. LightGBM o XGBoost
son la elección natural por su rendimiento en datos tabulares. Dado que 50.000 observaciones entrenan
muy rápido, se programa una búsqueda bayesiana de hiperparámetros para que corra de forma desatendida
(de tal manera que pueda uno descansar). Para evitar el sobreajuste, se restringe el espacio de búsqueda
a los parámetros críticos (regularización, submuestreo, profundidad) y se emplea \textit{early stopping}.

\medskip
Entre las horas dieciséis y dieciocho se construye un ensamblaje ligero, combinando las predicciones
del modelo de \textit{boosting} con las del modelo lineal de referencia mediante un promedio ponderado.
En datasets de tamaño moderado, combinar modelos complementarios suele aportar entre 0{,}5 y 2 puntos
de AUC sin un coste excesivo.

\medskip
Las horas dieciocho a veintidós se reservan para la calibración y la evaluación final: calibrar las
probabilidades si se requiere precisión probabilística; evaluar en el conjunto de test con métricas
alineadas al problema (AUC-ROC, AUC-PR, F1 en el umbral operativo); y verificar la consistencia entre
particiones. Las dos últimas horas se dedican a documentar las decisiones, los resultados y las
limitaciones.

\medskip
Las decisiones clave se justifican por la restricción temporal y operativa: el orden secuencial
(variables antes que hiperparámetros) evita rehacer trabajo; programar la búsqueda bayesiana de forma
desatendida optimiza el ciclo de 24 horas reales; y la validación estratificada en 5 particiones
aporta la robustez necesaria para un dataset de 50.000 observaciones frente a un simple \textit{hold-out}.
}

%% ─────────────────────────────────────────────
\item \textbf{¿Dónde se concentra el valor diferencial de la experiencia profesional?}

Una vez dominadas las herramientas técnicas (algoritmos, métricas, validación, despliegue),
¿en qué aspectos del trabajo un profesional experimentado aporta más valor diferencial frente a un
profesional competente pero con menor trayectoria? ¿Qué conocimiento se adquiere por experiencia
y no por formación?

\sol{

El profesional experimentado no se distingue por conocer más algoritmos, sino por un conjunto de
capacidades que solo emergen de la práctica repetida ante situaciones imprevisibles.

\medskip
La primera es el diagnóstico rápido de problemas: reconoce patrones de fallo (data leakage, fuga del
objetivo, cambios en la distribución de los datos) antes de entrenar el primer modelo, porque los ha
visto muchas veces, y no necesita experimentos para sospechar el diagnóstico correcto. La segunda
es la definición del problema adecuado: sabe que una de las mayores fuentes de fracaso no es el
modelo, sino la definición del objetivo, la elección de la ventana temporal o la confusión entre lo
que se puede medir y lo que de verdad importa para el negocio, un juicio que no suele enseñarse en un
curso. A ello se añade la calibración del esfuerzo, es decir, saber cuándo vale la pena seguir
optimizando y cuándo el rendimiento adicional ya no se traduce en valor de negocio; un perfil junior
tiende a optimizar hasta el límite técnico, mientras que un experto sabe cuándo detenerse.

\medskip
Igual de importante es la gestión de la incertidumbre y la comunicación: transmitir con precisión qué
puede y qué no puede garantizar el modelo, y traducir esa incertidumbre a términos que los
interesados no técnicos puedan usar para decidir. A esto se suma el conocimiento de los fallos en
producción, pues el experto ha comprobado que un modelo puede funcionar bien en desarrollo y fallar
en producción por causas invisibles durante este (skew entrenamiento-servicio, dependencias de datos,
degradación temporal), y ha interiorizado cómo prevenirlas. También cuenta la intuición sobre los
datos, la capacidad de detectar anomalías o inconsistencias que no se aprecian en las estadísticas
resumen, resultado acumulado de haber trabajado con muchos conjuntos de datos. Por último, la gestión
de los interesados: saber cuándo decir que no, cómo negociar plazos realistas y cómo convertir un
proyecto técnico en un proyecto de negocio con impacto demostrable.
}

%% ─────────────────────────────────────────────
\item \textbf{Discrepancia entre entorno de desarrollo y producción.}

Enumere las causas más frecuentes por las que un modelo con buenas métricas en desarrollo falla
al desplegarse. Para cada causa, explique qué la hace difícil de detectar y cómo prevenirla desde
el diseño.

\sol{

\begin{enumerate}[1.]
  \item El training-serving skew se produce cuando el modelo se entrena con datos procesados de una
  manera y recibe en producción datos procesados de forma ligeramente distinta (otro orden de
  operaciones, otras versiones de librerías, otro tratamiento de los nulos). Es difícil de detectar
  porque las diferencias son sutiles y el modelo no falla con un error explícito, sino que
  simplemente se degrada. Se previene usando el mismo pipeline de transformación, con idéntico código
  y artefacto serializado, en entrenamiento e inferencia.

  \item El data leakage en desarrollo aparece cuando el modelo usa variables que en desarrollo
  parecen inocuas pero que en producción no están disponibles en el momento de la predicción
  (información futura o variables calculadas tras el evento). Resulta difícil de detectar porque las
  métricas son excelentes y la fuga es conceptual, no técnica. Se previene documentando el instante de
  disponibilidad de cada variable y auditando la posible fuga antes de entrenar.

  \item El data drift consiste en que la distribución de los datos en producción difiere de
  la del entrenamiento, ya sea por estacionalidad, por cambios en el comportamiento de los usuarios o
  por cambios en los sistemas de captura. Es difícil de detectar durante el desarrollo porque el
  conjunto de test es estático. Se previene monitorizando en producción la distribución de las
  variables de entrada (PSI, divergencia KL o el test de Kolmogorov-Smirnov) y estableciendo alertas.

  \item Las dependencias frágiles de datos surgen cuando el modelo depende de una fuente que en
  producción llega con retraso, con errores ocasionales o con un esquema cambiante. Es difícil de
  detectar porque en desarrollo los datos suelen estar completos y limpios. Se previene diseñando el
  pipeline con un manejo explícito de los datos faltantes, validación del esquema de entrada y cortes
  automáticos ante datos anómalos.

  \item La inconsistencia de versiones se debe a diferencias en las versiones de las librerías
  (scikit-learn, XGBoost) entre desarrollo y producción, que pueden producir predicciones
  numéricamente distintas. Es difícil de detectar porque el modelo no falla, sino que ofrece
  resultados levemente diferentes que degradan el rendimiento de forma imperceptible. Se previene
  fijando las versiones en un fichero de requisitos y trabajando en un entorno reproducible.

  \item La latencia y las restricciones de cómputo aparecen cuando el modelo requiere más tiempo de
  inferencia del aceptable en producción, lo que obliga a simplificarlo tras el despliegue de forma no
  controlada. Se previene definiendo los límites de latencia antes de entrenar y midiendo el tiempo de
  inferencia como parte del proceso de evaluación.
\end{enumerate}
}

%% ─────────────────────────────────────────────
\item \textbf{Discrepancia entre métrica y utilidad operativa.}

Un modelo de clasificación obtiene un AUC de 0{,}95 en el conjunto de test, pero los usuarios del
sistema reportan que las predicciones no son útiles en la práctica. Explique cómo es posible esta
discrepancia. ¿Qué diagnóstico realizaría?

\sol{

El AUC-ROC mide la capacidad discriminativa del modelo en media sobre todos los umbrales posibles. Un
AUC de 0{,}95 garantiza que el modelo ordena bien los ejemplos positivos y negativos, pero no dice
nada sobre su utilidad en el umbral operativo concreto ni sobre si las predicciones responden a la
pregunta que los usuarios realmente necesitan resolver. Esa brecha entre una métrica agregada y la
decisión concreta explica la discrepancia, y admite varias causas.

\medskip
Puede ocurrir que el objetivo sea incorrecto, de modo que el modelo prediga un sustituto del evento
de interés y no el evento en sí: las métricas son excelentes, pero las predicciones no se
corresponden con las decisiones que deben tomarse. También es habitual una calibración deficiente,
en la que el modelo discrimina bien pero las probabilidades predichas están mal ajustadas; un modelo
que predice casi siempre 0{,}99 o 0{,}01 puede tener un AUC alto y resultar inútil para cualquier
decisión basada en umbrales de probabilidad. En problemas muy desbalanceados (muchos datos de una clase y pocos de la otra), el AUC puede
ser engañoso, pues, por ejemplo, para un problema con la mayor parte de las observaciones reales negativas, un modelo que clasifica todo como negativo alcanza un valor alto por la forma en
que se construye la curva ROC; en ese caso el AUC-PR (área bajo la curva precisión-recall) suele ser
más informativo. Otra posibilidad es que el umbral operativo caiga en una zona de bajo rendimiento:
el AUC es alto en promedio, pero en el rango de probabilidades donde se toman las decisiones (por
ejemplo, entre 0{,}4 y 0{,}6) el modelo discrimina mal. A ello se añaden una distribución del
conjunto de test no representativa, construida con un sesgo de selección que no refleja la realidad de
producción, y la desalineación entre la métrica y el objetivo de negocio, ya que el AUC optimiza el
ordenamiento, pero el negocio puede necesitar precisión alta en el tramo superior de las
predicciones, maximizar el recall o minimizar los falsos positivos en un contexto de coste asimétrico.

\medskip
El diagnóstico pasa por analizar la curva de calibración, el AUC-PR y la matriz de confusión en el
umbral operativo real, validar la distribución del conjunto de test frente a la de producción y,
sobre todo, entrevistar a los usuarios para entender qué tipo de predicción esperan y cómo la
utilizan.
}

%% ─────────────────────────────────────────────
\item \textbf{Detección de regímenes de mercado.}

Supón que deseas identificar automáticamente distintos regímenes de mercado (alcista, bajista, alta
volatilidad, crisis, recuperación, etc.) utilizando información histórica de rentabilidades,
volatilidades y factores macroeconómicos. Diseña una metodología basada en reducción dimensional y
clustering para abordar este problema. Justifica qué algoritmos elegirías y cómo validarías que los
regímenes identificados tienen significado económico. Explica también las limitaciones que podrían
surgir si los datos contienen ruido, observaciones atípicas o estructuras no esféricas.

\sol{

El estado del mercado en cada periodo se caracteriza mediante una ventana temporal rodante que reúne
las rentabilidades realizadas a distintos horizontes (1, 5 y 21 días), la volatilidad realizada, la
correlación entre activos, los spreads de crédito, los tipos de interés, indicadores
como el VIX y factores macroeconómicos. Cada observación es así un vector que representa el estado del
mercado en ese instante.

\medskip
Con entre 20 y 50 variables, se aplica PCA para extraer entre tres y cinco componentes principales que
resuman ese estado; conviene tratar antes los valores atípicos, ya que el PCA es sensible a ellos, y
para la exploración visual puede recurrirse a UMAP en dos dimensiones.

\medskip
Sobre esa representación reducida, el clustering se aborda con varios algoritmos. K-Means
(con $k$ entre 4 y 6) es rápido y reproducible y permite interpretar los centroides como perfiles típicos de
cada régimen, por lo que constituye una primera opción razonable por su sencillez. Los Hidden Markov
Models resultan especialmente adecuados para series temporales con dependencia secuencial, ya que los
regímenes presentan persistencia y transiciones con probabilidades estimadas, lo que los hace más
apropiados que un K-Means puro para datos financieros con inercia. Los Gaussian Mixture Models
generalizan K-Means a clústeres elípticos y permiten una asignación probabilística (cada periodo
pertenece a varios regímenes con cierta probabilidad), útil en las transiciones entre regímenes.

\medskip
La validación del significado económico se apoya en tres comprobaciones: que la distribución de
rentabilidades del índice de referencia difiera de forma apreciable entre regímenes; que los
regímenes identificados coincidan con periodos económicos reconocidos (2008, 2020 y 2022); y que
las estrategias de asignación de activos condicionadas al régimen mejoren el Sharpe respecto a una
estrategia estática.

\medskip
Las limitaciones más relevantes proceden de la naturaleza de los datos. El ruido de alta frecuencia
es elevado, y aunque la reducción dimensional ayuda, no lo elimina, de modo que las ventanas
temporales deben ser lo bastante largas para suavizarlo sin perder la información de los cambios de
régimen. Los outliers son inevitables, pues las crisis son por definición observaciones atípicas que
pueden dominar el PCA, lo que aconseja tratarlos previamente. Las estructuras no esféricas también
suponen un problema: los regímenes no son necesariamente convexos ni esféricos, y mientras K-Means
asume esfericidad, los Gaussian Mixture Models admiten elipses y DBSCAN puede captar formas
arbitrarias a costa de una menor estabilidad. Por último, hay que evitar la fuga de información del
futuro: si los regímenes se definen con el histórico completo y luego se usan como variables
predictivas, se introduce un sesgo, por lo que cualquier estrategia basada en regímenes debe
definirlos usando solo los datos disponibles hasta el momento de la predicción.
}

%% ─────────────────────────────────────────────
\item \textbf{Selección de un único algoritmo para datos tabulares: argumentación completa.}

Si solo pudiera utilizar un algoritmo de ML para datos tabulares durante toda su carrera
profesional, ¿cuál elegiría? Argumente considerando: versatilidad, robustez ante datos imperfectos,
velocidad de iteración, competitividad como baseline y facilidad de despliegue.

\sol{

La elección sería LightGBM, un algoritmo de árboles potenciados por gradiente, y los cinco criterios
propuestos lo respaldan.

\medskip
En cuanto a la versatilidad, LightGBM resuelve clasificación binaria, multiclase y regresión con el
mismo algoritmo y cambios mínimos de configuración, y trabaja con variables numéricas, categóricas
(con codificación nativa) y datos con valores faltantes sin un preprocesamiento especial, sin
necesidad siquiera de normalizar. Respecto a la robustez ante datos imperfectos, los árboles de
decisión son por naturaleza poco sensibles a los outliers, a las variables irrelevantes y a las
distribuciones no normales, y la regularización integrada ($L_1$, $L_2$) ayuda a controlar el
sobreajuste incluso con muestras pequeñas. En velocidad de iteración, LightGBM tiende a ser más
rápido que XGBoost estándar gracias al crecimiento de árboles por hojas (leaf-wise) y a otras
optimizaciones de eficiencia, lo que acorta los ciclos de experimentación, algo especialmente valioso
cuando hay restricciones de tiempo.

\medskip
Sobre su competitividad como referencia, LightGBM con hiperparámetros por defecto suele rendir muy
bien en datos tabulares estructurados y destaca con frecuencia tanto en competiciones como en
comparativas frente a las redes neuronales tabulares. Y en cuanto a la facilidad de despliegue, el
modelo serializado es un artefacto ligero, con una inferencia rápida y de latencia predecible que no
depende de GPU.

\medskip
La razón para preferirlo como única apuesta de por vida no es que sea siempre el mejor, sino que rara
vez es una mala elección: cuando no gana, suele quedar cerca, y su ecosistema de herramientas (SHAP
nativo, early stopping, validación cruzada integrada) lo hace cómodo de usar en casi cualquier
problema tabular.
}

%% ─────────────────────────────────────────────
\item \textbf{Contextos donde el deep learning no es la solución apropiada.}

Describe al menos cuatro situaciones concretas donde una red neuronal es una elección inadecuada
frente a alternativas más simples. Justifique las razones de fondo en cada caso.

\sol{

\begin{enumerate}[1.]
  \item Datos tabulares estructurados con pocas observaciones ($n < 10.000$). Las redes neuronales
  suelen necesitar grandes volúmenes de datos para superar el sobreajuste, de modo que con muestras
  pequeñas los árboles de decisión regularizados (XGBoost, LightGBM) o los modelos lineales
  regularizados tienden a rendir mejor. La razón de fondo es que una red neuronal puede tener un
  número de parámetros muy elevado, difícil de estimar de forma fiable con pocos datos.

  \item Entornos con requisitos estrictos de interpretabilidad regulatoria. En la concesión de
  crédito, el diagnóstico médico asistido o las decisiones judiciales, la regulación exige explicar
  cada predicción individual. Los modelos lineales y los árboles de decisión producen explicaciones
  exactas y verificables, mientras que las redes neuronales solo ofrecen aproximaciones posteriores
  (SHAP, LIME), difíciles de auditar con el mismo rigor.

  \item Tiempo de cómputo o recursos limitados. Entrenar una red neuronal competitiva suele requerir
  GPU, muchas iteraciones de entrenamiento y una búsqueda intensiva de hiperparámetros (arquitectura,
  tasa de aprendizaje, dropout, normalización por lotes). Para un ciclo de iteración rápido o un
  despliegue en hardware modesto, un modelo de gradient boosting entrenado en minutos es preferible.

  \item Señal débil en datos tabulares de alta dimensión y mucho ruido. Las redes neuronales pueden
  amplificar el ruido cuando la relación señal-ruido es baja, porque aprenden con facilidad patrones
  espurios. Los modelos con regularización explícita (Lasso, Ridge, o árboles con la profundidad y el
  tamaño de hoja limitados) son más conservadores y suelen generalizar mejor en estos contextos.

  \item Problemas con variables muy interpretables y efectos aditivos conocidos. Si el conocimiento de
  dominio indica que la relación entre las variables y el objetivo es aproximadamente lineal o
  log-lineal, una regresión logística o una regresión lineal regularizada es más eficiente, más
  interpretable y menos propensa al sobreajuste que una red neuronal que dedicaría capacidad a modelar
  no linealidades inexistentes.
\end{enumerate}
}

%% ─────────────────────────────────────────────
\item \textbf{Criterios para reentrenar un modelo desplegado.}

Un modelo fue desplegado hace seis meses. ¿Qué indicadores monitoriza para determinar si necesita
reentrenamiento? ¿Cómo distingue un deterioro real de una fluctuación normal?

\sol{

Para monitorizarlo, podemos utilizar varios indicadores complementarios. El más directo es la métrica de
rendimiento en producción: si se dispone de labels con una latencia aceptable (por ejemplo,
cuando el evento ocurre en los 30 días siguientes), conviene seguir la métrica principal (AUC, F1,
error cuadrático) en ventanas rodantes, de modo que un deterioro sostenido por encima de un umbral
predefinido sea señal de reentrenamiento. A ello se añade la vigilancia del data drift, comparando la
distribución de las variables de entrada en producción con la del entrenamiento mediante el PSI
(índice de estabilidad poblacional, que mide cuánto se ha desplazado una variable) o el test de
Kolmogorov-Smirnov; un PSI elevado en variables críticas suele tomarse como indicio de un drift
relevante. También conviene vigilar el concept drift, es decir, un cambio acusado en la distribución
de las predicciones sin que cambien las entradas, que puede revelar una alteración en la relación
entre variables y objetivo, así como la tasa de positivos observados, ya que una divergencia marcada
frente a la tasa predicha apunta a una pérdida de calibración. En última instancia, el deterioro
relevante es el que se refleja en las métricas de negocio (conversión, pérdidas, retención).

\medskip
Distinguir un deterioro real de una fluctuación normal exige cierta disciplina estadística. Resulta
útil emplear bandas de control (media $\pm 2\sigma$ o $\pm 3\sigma$ sobre la historia reciente de la
métrica) y activar la alerta solo cuando la métrica cae fuera de la banda de manera persistente, no
puntual. Conviene además exigir que el deterioro se mantenga durante un mínimo de periodos
consecutivos (por ejemplo, tres semanas) antes de reentrenar, para no reaccionar a fluctuaciones
aleatorias, y comparar el periodo deteriorado con periodos históricos de estacionalidad similar, dado
que un descenso que coincide con un patrón estacional conocido no requiere reentrenamiento sino un
ajuste del umbral operativo. Por último, si el cambio en la distribución de las entradas es abrupto y
no gradual, la causa puede ser un error de datos en una fase anterior del pipeline, y no un deterioro
del modelo, por lo que conviene investigar la fuente antes de reentrenar.
}

%% ─────────────────────────────────────────────
\item \textbf{El tradeoff sesgo-varianza en las decisiones cotidianas.}

El tradeoff sesgo-varianza es el concepto teórico más citado del ML. Sin embargo, rara vez se
invoca explícitamente en la práctica diaria. Explique cómo este tradeoff se manifiesta de forma
implícita en las decisiones cotidianas que un profesional de ML toma sin nombrar el concepto.
Proporcione al menos tres ejemplos concretos.

\sol{

El tradeoff sesgo-varianza establece que el error esperado de un modelo se descompone en sesgo
(la incapacidad estructural del modelo para capturar la relación real), varianza (su sensibilidad a
las fluctuaciones del conjunto de entrenamiento) y un término irreducible. Sus implicaciones están presentes en numerosas
decisiones.

\begin{enumerate}[1.]
  \item Al elegir el número de árboles y la profundidad en XGBoost. Cuando se sube la profundidad
  máxima de 4 a 8 se reduce el sesgo, porque el modelo puede capturar relaciones más complejas, a
  costa de aumentar la varianza, ya que se vuelve más sensible a los datos de entrenamiento. Y al usar
  early stopping para detener el entrenamiento antes de que la métrica de validación empeore, se está
  controlando la varianza.

  \item Al decidir si se aplica regularización Lasso. Al preferir Lasso a la
  regresión ordinaria en un modelo con 50 variables y 1.000 observaciones, se está introduciendo sesgo
  (al asumir que muchos coeficientes son cero) para reducir la varianza.

  \item Al elegir entre una regresión lineal y un GBM para un dataset pequeño. Cuando, con 500
  observaciones, se decide empezar por un modelo simple y solo aumentar la complejidad si
  hace falta, se está priorizando una varianza baja sobre un sesgo bajo. El modelo lineal tendrá sesgo
  alto si la relación es no lineal, pero con tan pocos datos la varianza de un GBM probablemente
  dominaría el error total.

  \item Al fijar el tamaño de la ventana de entrenamiento en series temporales. Una ventana corta
  produce estimaciones más adaptadas al régimen actual (bajo sesgo temporal) pero más ruidosas (alta
  varianza), mientras que una ventana larga es más estable pero puede incorporar regímenes obsoletos.
  Esta decisión es en sí gestión del tradeoff sesgo-varianza en el tiempo.
\end{enumerate}
}

%% ─────────────────────────────────────────────
\item \textbf{Feature engineering: ¿cuándo es necesario y cuándo no?}

Discuta en qué contextos la construcción manual de variables sigue siendo determinante para el
rendimiento del modelo y en qué contextos no aporta valor significativo. Justifique la diferencia.

\sol{

La construcción manual de variables es determinante, ante todo, cuando el conocimiento del dominio no
está presente en los datos en bruto. En finanzas, por ejemplo, un ratio como la deuda neta en
relación con los beneficios resume en una sola cifra la solidez de una empresa, algo que un modelo
alimentado solo con cifras contables sueltas necesitaría una enorme cantidad de datos para deducir por
sí mismo. Ocurre algo parecido en las series temporales con estacionalidad: extraer de forma explícita
el día de la semana, los festivos o indicadores de los ciclos que se repiten ayuda a que los modelos
sin memoria del orden temporal (los árboles o los modelos lineales) reconozcan esos patrones. Y cuando
hay mucho ruido y poca señal, agregar la información (por ejemplo, con medias móviles o variaciones
porcentuales) suaviza el ruido y resalta lo relevante, sobre todo cuando hay pocos datos. Un último
caso son las variables categóricas con muchísimos valores distintos: en lugar de crear una columna por
valor (lo que se conoce como one-hot), suele funcionar mejor resumir cada categoría con un número, como
su frecuencia de aparición o su relación media con el objetivo.

\medskip
En cambio, construir variables a mano aporta poco cuando el propio modelo ya es capaz de extraer esa
información. Con imágenes, texto o audio, las redes neuronales profundas (como las redes
convolucionales) aprenden por sí solas qué conviene representar, y una transformación hecha a mano
(por ejemplo, un histograma de colores de una imagen) se queda corta frente a lo que aprende la red.
Tampoco compensa demasiado con conjuntos de datos grandes y de buena calidad, porque ahí los modelos
de boosting ya captan solos las interacciones y las relaciones no lineales. Y si el tiempo de
desarrollo es escaso, mantener y validar variables complicadas en producción puede costar más de lo
que aportan.
}

%% ─────────────────────────────────────────────
\item \textbf{Construcción de un modelo de ensamblaje para predicción de retornos financieros.}

Suponga que se trabaja en el equipo cuantitativo de un Hedge Fund y se dispone de una batería de
modelos desarrollados por distintos investigadores para predecir la rentabilidad futura de las
acciones. Algunos modelos están basados en variables fundamentales, otros en factores técnicos, otros
en datos macroeconómicos y otros en variables alternativas. Explique cómo diseñaría una estrategia
de ensamblaje que combine las predicciones individuales en un único modelo global. Discuta qué tipos
de modelos utilizar como modelos base, cómo evaluar el rendimiento individual de cada uno y qué
técnicas de ensamblaje emplear para combinar sus predicciones. Analice, además, cómo determinar si
el ensamblaje aporta valor añadido respecto a los modelos individuales y qué riesgos de sobreajuste
podrían surgir durante el proceso. El objetivo no es encontrar el mejor modelo, sino construir un
sistema donde los errores de unos modelos sean compensados por las fortalezas de otros.

\sol{

El valor del ensamblaje proviene de la diversidad, no de la calidad individual: dos modelos mediocres
pero con errores poco correlacionados pueden generar, al combinarse, un predictor mejor que un único
modelo más preciso. El objetivo, por tanto, es maximizar la diversidad de las fuentes de error. Como
modelos base conviene partir de familias complementarias, en la línea del resto del pipeline: modelos
lineales regularizados (Ridge o Lasso), modelos de boosting (XGBoost o LightGBM) y algún modelo de
factores alineado con la lógica financiera, de manera que cada uno capte una parte distinta de la
señal.

\medskip
La evaluación individual de los modelos base debe ir más allá de su acierto medio. Para cada modelo
conviene medir, en validación walk-forward, la correlación de Spearman entre sus predicciones y la
rentabilidad realizada, así como la consistencia de esa correlación a lo largo del tiempo (su media
frente a su desviación), pues una señal muy volátil delata un modelo dependiente del régimen. La
métrica clave para la diversidad, sin embargo, es la correlación de los errores entre modelos:
aquellos pares con baja correlación de errores son los candidatos ideales para combinarse. Conviene
además vigilar si la capacidad predictiva de cada modelo se deteriora con el tiempo, lo que obligaría
a reentrenarlo con más frecuencia.

\medskip
Para combinar las predicciones, una opción sencilla y robusta es el promedio simple, que resiste bien
el sobreajuste y, aunque existen métodos de combinación más elaborados, suele ser difícil de batir
fuera de muestra cuando no hay evidencia sólida de que un modelo sea sistemáticamente superior. Otra
opción, la misma que se emplea en el pipeline, es el promedio ponderado, que asigna más peso a los
modelos con mejor rendimiento reciente y permite así adaptarse a los cambios de régimen; conviene
mantener esos pesos estables y rebajar el de los modelos muy redundantes, para no concentrar la
combinación en una única fuente de señal.

\medskip
El ensamblaje aporta valor añadido si, fuera de muestra, su correlación con el retorno y su
consistencia superan a las del mejor modelo individual, y si esa mejora se traslada a la estrategia:
mayor Sharpe y menor drawdown en una cartera long/short que con cualquier modelo por separado.

\medskip
Los riesgos de sobreajuste en este proceso son específicos. Si los pesos se ajustan mirando todo el
histórico, se acaban capitalizando correlaciones pasajeras, por lo que conviene fijarlos con
validación walk-forward y mantener la combinación sencilla. Incorporar muchos modelos muy
correlacionados aporta poca diversidad real: un conjunto numeroso de modelos que comparten la mayoría
de las variables tiende a aportar menos que unos pocos con fuentes de información distintas. Y, sobre
todo, los pesos deben estimarse sin usar datos posteriores a la fecha de predicción, para no
introducir una fuga de información.
}

%% ─────────────────────────────────────────────
\item \textbf{Tratamiento de la clasificación con clases desbalanceadas.}

En un problema con un 98\% de negativos y un 2\% de positivos, un compañero aplica SMOTE y reporta
que ``el accuracy ha mejorado''. Analice críticamente esta afirmación y describa un enfoque más
adecuado.

\sol{

La afirmación es muy probablemente irrelevante. En un dataset con un 98\,\% de negativos, un modelo
trivial que clasifica todo como negativo ya alcanza un accuracy del 98\,\%, de modo que cualquier
mejora de accuracy reportada tras aplicar SMOTE debe compararse con esa referencia. Lo más probable es
que SMOTE haya aumentado el recall de la clase positiva (lo cual sería deseable) a costa de
incrementar los falsos positivos, con un efecto neto ambiguo que el accuracy no captura.

\medskip
SMOTE plantea además problemas específicos en este contexto. Genera ejemplos sintéticos interpolando
entre positivos reales en el espacio de variables, de manera que si los positivos son heterogéneos o
forman grupos dispersos, esos ejemplos pueden introducir ruido en regiones de alta densidad negativa.
Si se aplica antes de la validación cruzada, en lugar de dentro de cada partición, produce data
leakage, porque ejemplos sintéticos basados en el conjunto de validación contaminan el entrenamiento
e inflan artificialmente las métricas. Y, en cualquier caso, el accuracy es una métrica inadecuada
para clases desbalanceadas, pues no penaliza los errores en la clase minoritaria, que es precisamente
la que importa.

\medskip
Un enfoque más adecuado seguiría estos pasos:

\begin{enumerate}[1.]
  \item Elegir la métrica correcta: AUC-PR (área bajo la curva precisión-recall), F1 de la clase
  positiva o una función de pérdida ponderada alineada con el coste asimétrico del error (falso
  negativo frente a falso positivo).
  \item Ponderar la pérdida por clase, ya que muchos algoritmos permiten dar más peso a los errores
  sobre la clase minoritaria sin necesidad de generar datos sintéticos.
  \item Ajustar el umbral de clasificación, desplazándolo por debajo del 0{,}5 por defecto para
  aumentar el recall según el coste de negocio.
  \item Recurrir a estrategias de muestreo con cuidado: si se utiliza SMOTE, aplicarlo dentro de cada
  partición de la validación cruzada y nunca sobre el dataset completo antes de dividir.
  \item Valorar el submuestreo de la clase mayoritaria, que en datasets muy grandes puede ser más
  eficiente y menos ruidoso que la generación sintética de positivos.
\end{enumerate}
}

%% ─────────────────────────────────────────────
\item \textbf{Cuantificación de la incertidumbre: ¿cuándo es prescindible y cuándo es crítica?}

Un modelo probabilístico (que produce intervalos de confianza) es más complejo de construir y
mantener que uno que produce solo predicciones puntuales. Describe dos contextos donde la
incertidumbre es prescindible y dos donde su ausencia supone un riesgo grave. Explica la diferencia
de fondo.

\sol{

La incertidumbre es prescindible, en primer lugar, en la recomendación de contenido a escala. Un
sistema que ordena 1.000 películas para cada usuario no necesita intervalos de confianza: la decisión
de mostrar los diez primeros resultados es de bajo coste, reversible al recargar la página, y el
propio usuario evalúa el resultado de inmediato, de modo que el volumen y la reversibilidad eliminan
la necesidad de cuantificar la incertidumbre. Algo parecido ocurre en la clasificación de spam en
sistemas de alto volumen, ya que el objetivo del sistema es facilitar al usuario la tarea de revisar su bandeja de entrada, éste, no desea obtener para cada correo las probabilidades de ser spam y revisarlos uno a uno, es decir, una predicción puntual con un umbral bien calibrado basta, y la complejidad añadida de un modelo probabilístico no se traduce
en valor.

\medskip
Su ausencia, en cambio, supone un riesgo grave cuando la decisión es irreversible o el coste del error
es asimétrico. En la gestión de riesgo financiero, estimar la magnitud de una posible pérdida sin
acompañarla de una banda de incertidumbre puede inducir a asumir una exposición excesiva que ignore la
cola de la distribución de pérdidas; aquí un modelo probabilístico, que aporta esa incertidumbre
calibrada, es muy preferible a una predicción puntual. Lo mismo sucede en el diagnóstico médico
asistido: un modelo que estima la probabilidad de una enfermedad grave sin comunicar su incertidumbre
puede llevar a un médico a una decisión clínica que él mismo rechazaría si supiera que la predicción
tiene una varianza enorme, pues la ausencia de incertidumbre crea una falsa sensación de precisión; en
decisiones irreversibles y de alto impacto (cirugía, quimioterapia), la incertidumbre forma parte de
la información clínicamente relevante.

\medskip
Es decir, la incertidumbre resulta prescindible cuando la decisión es reversible, el coste del
error es bajo, y el volumen de decisiones permite que los errores se promedien; y se
vuelve crítica cuando la decisión es irreversible, el coste del error es muy alto, y
quien decide necesita saber cuánta confianza depositar en la predicción para actuar de forma
responsable.
}

\end{enumerate}

\noindent \textit{Nota:} Se ha utilizado Claude Opus 4.7. como asistente para la resolución de los ejercicios y la escritura en \LaTeX.

\end{document}
